\SourceLeaf{146}
senden, doch erlauben Sie wohl, dass ich es erst noch einige Wochen behalte. Bei vorläufigem Durchblättern scheint es mir sehr viel Schönes zu enthalten. Vieles von dem, was ich in meiner Methode besonders in ihrer ersten Gestalt Eigenthümliches hatte, finde ich auch in diesem Buche wieder. Es scheint mein Schicksal zu sein, fast in allen meinen theoretischen Arbeiten mit \textsc{Legendre} zu concurriren. So in der höhern Arithmetik, in den Untersuchungen über transcendente Functionen, die mit der Rectification der Ellipse zusammenhangen, bei den ersten Gründen der Geometrie und nun wieder hier. So ist z.\,B.\ auch das von mir seit 1794 gebrauchte Princip, dass man, um mehrere Grössen, die man nicht alle genau darstellen kann, am besten darzustellen, die Summe der Quadrate zu einem Minimum machen müsse, auch in \textsc{Legendres} Werke gebraucht und recht wacker ausgeführt.

\begin{center}
. . . . . . . . . . . . .
\end{center}

\begin{center}
\textsc{Gauss} an \textsc{Olbers}. Göttingen, 4. October 1809.
\end{center}

. . . . . . . . . Erinnern Sie sich wohl noch, liebster Freund, dass ich bei meiner ersten Anwesenheit in Bremen 1803 mit Ihnen über das Princip gesprochen habe, dessen ich mich bediente, Beobachtungen am genauesten darzustellen, dass nemlich bei gleichem Werthe der Beobachtungen die Summe der Quadrate der Differenzen ein Kleinstes sein muss? Dass wir darüber 1804 in Rehburg gesprochen haben, davon sind mir noch alle Umstände gegenwärtig. Es ist mir daran gelegen, diess zu wissen. Über die Ursache der Frage ein andermal.

\begin{center}
. . . . . . . . . . . . . . .
\end{center}
